% bab3.tex — dari BAB_III_PENUTUP_POLISHED.md mulai ## A.

\phantomsection
\section*{BAB III --- PENUTUP}
\addcontentsline{toc}{section}{BAB III --- PENUTUP}

\phantomsection
\subsection*{A. Kesimpulan}
\addcontentsline{toc}{subsection}{A. Kesimpulan}

Pertama, taksonomi hadis merumuskan lima tahapan: pembentukan multi-saluran kenabian, penyebaran sahabat, pembakuan transmisi tabi'in, pembukuan kanonik ulama klasik, dan keragaman metodologis modern.\footnote{T. N. S. Tanjung and Z. Fadhlurrahman, "The role of tadwin hadith in the development of Islamic historiography," \\textit{Istifham: Journal of Islamic Studies}, vol. 3, no. 3, hlm. 188--201, 2025; A. Musadad, "Articulating the prophetic authority: Some notes on the nature of hadith collection in al-Muwatta' and musannaf literature," \\textit{Jurnal Studi Ilmu-Ilmu Al-Qur'an dan Hadis}, vol. 27, no. 1, hlm. 351--373, 2026.} Pascakeguncangan \textit{al-fitnah al-kubra}, keharusan \textit{sanad} dan semangat \textit{tathabbut} dipaksakan oleh keadaan; \textit{rihlah} kemudian mengukuhkannya sebagai kelaziman teknis hingga program \textit{tadwin} masa Umar bin Abdul Aziz.\footnote{Tanjung, "The Role of Tadwin Hadith."; N. M. A. M. Razak et al., "Evolution of hadith textual criticism: From classical to contemporary approaches," \\textit{Global Journal Al-Thaqafah}, hlm. 15--29, 2025; Musadad, "Articulating the Prophetic Authority.".} Berikutnya, penghimpunan \textit{Kutub al-Sittah} mematangkan kodifikasi, sedangkan trikotomi sahih-\textit{hasan}-\textit{dhaif} diberlakukan sebagai ukuran keotentikan berjenjang yang mengikat sarjana hingga kini. Rumusan masalah kesatu pun terjawab: tiap zaman memformalkan semangat penjagaan pendahulu ke perangkat yang kian baku --- dari ingatan-terapan ke kanon dan pangkalan data digital.

Kedua, \textit{hifz} ditegakkan sebagai basis kultural, \textit{kitabah} menopang ingatan, hingga \textit{tadwin} bersponsor penguasa menjadi muaranya.\footnote{H. Rashwan, "Hadith as oral literature through early Islamic literary criticism," \\textit{Studia Islamica}, vol. 119, no. 1, hlm. 34--69, 2024; M. Lutfianto et al., "History of hadith writing in the precodification period in the prospective of 'Ajjaj al-Khatib," \\textit{Al-Thiqah: Jurnal Ilmu Keislaman}, vol. 7, no. 2, hlm. 308, 2024.} Lewat jejaring guru-murid pengendali \textit{hifz}, \textit{sanad} berlandaskan \textit{'adalah} dan \textit{dabt} pun dilahirkan; sengketa pelarangan-perizinan ditengahi --- dari \textit{sadd al-dzari'ah} ke \textit{rukhsah}\footnote{H. B. Jaiyeoba and N. M. Osmani, "Hadith preservation: Techniques and contemporary efforts," \\textit{Journal of Fatwa Management and Research}, vol. 29, no. 3, hlm. 31--45, 2024; S. Watukila, "Historiografi hadis di masa mutaqaddimin dalam pandangan Mustafa Azami," \\textit{Amsal Al-Qur'an: Jurnal Al-Qur'an dan Hadis}, vol. 1, no. 3, hlm. 221--235, 2024.} --- sebagaimana dibuktikan \textit{Sahifah al-Sadiqah}. Penyebaran sahabat ke pelbagai wilayah menghasilkan khazanah kedaerahan yang tersaring lewat \textit{tathabbut} Khulafaur Rasyidin, lalu dirapikan melalui program \textit{tadwin} yang berpuncak pada \textit{al-Muwatta'} dan pola kanonik al-Bukhari sebagai kulminasinya.\footnote{R. E. Gil, "Tracing the rihla in buldan books: Ibn al-Faqih's muhaddith resources," \\textit{Hitit Theology Journal}, vol. 23, hlm. 140--154, 2024; S. M. S. Obeid and F. S. Hussein, "Al-Bukhari's sources," \\textit{Jordan Journal for History and Archaeology}, vol. 17, no. 3, hlm. 44--62, 2023.} Terjawablah rumusan kedua: penghimpunan berlangsung bertahap-berlapis --- khazanah tulis menyangga kelisanan; prakarsa perseorangan ditata otoritas publik hingga menjelma kanon bersama.

Ketiga, empat gugus pendorong pemalsuan teridentifikasi --- politis, teologis-sektarian, destruktif-infiltratif, devosional sosio-ekonomis --- dilacak melalui pemeriksaan rangkap \textit{sanad}-\textit{matan} berpasangan.\footnote{Wasman et al., "A critical approach to prophetic traditions: Contextual criticism in understanding hadith," \\textit{Al-Jami'ah: Journal of Islamic Studies}, vol. 61, no. 1, hlm. 1--17, 2023; M. A. M. Muhamed Ali et al., "Al-jarh wa al-ta'dil (criticism and praise): Its significance in the science of hadith," \\textit{Mediterranean Journal of Social Sciences}, vol. 6, no. 2S1, hlm. 284, 2015.} Dorongan politis tumbuh dari kegentingan legitimasi pasca-\textit{fitnah} lewat \textit{afterlife} khilafah tiga dasawarsa; sementara unsur teologis, \textit{zindiq}, dongeng \textit{qussas}, dan anjuran \textit{targhib-tarhib} melanggengkan \textit{maudu'}.\footnote{H. H. Liew, "The caliphate will last for thirty years: Polemic and political thought in the afterlife of a prophetic hadith," \\textit{Journal of Islamic Studies}, vol. 36, no. 1, hlm. 38--82, 2025; I.-W. Su, "'Ali b. al-Madini (161-234/778-849): A critical reconstruction of his biography and evaluation of his contribution to hadith criticism," \\textit{Journal of Islamic Studies}, vol. 33, no. 1, hlm. 1--34, 2021; Wasman, "A Critical Approach to Prophetic Traditions."; U. M. Noor, "Muqaddimah of Ibn al-Salah and the revival of hadith studies in the Mamluk era," \\textit{Al-Bayan: Journal of Qur'an and Hadith Studies}, vol. 21, no. 2, hlm. 271--297, 2023.} Pendeteksian memadukan uji \textit{sanad} --- keutuhan periwayat, ketersambungan rantai, terbuktinya \textit{liqa'} --- dengan telaah \textit{matan}: keselarasan Qur'an, dukungan riwayat kukuh, kesesuaian historis; vonis dijatuhkan bila indikator\footnote{Muhamed Ali, "Al-jarh Wa Al-ta'dil (criticism and Praise)."; M. A. Calgan, "Analysis of Ibn Kathir's content criticism of some of the Sahihayn hadiths in the sirah section of al-Bidaya wa al-Nihaya," \\textit{Cumhuriyet Ilahiyat Dergisi}, vol. 28, no. 1, hlm. 303--324, 2024.} terpenuhi secara kumulatif-kontekstual. Dengan demikian, separuh rumusan ketiga terjawab: pemalsuan berwatak krisis otoritas struktural, sehingga pendeteksiannya mensyaratkan kritik \textit{sanad-cum-matan} yang terlembaga dan berkesinambungan.

Keempat, penyelamatan disusun dalam tiga lapis: \textit{takhrij}, imunitas \textit{jarh-ta'dil}, dan penghimpunan \textit{maudu'at} sebagai gudang memori penangkal disinformasi.\footnote{Muhamed Ali, "Al-jarh Wa Al-ta'dil (criticism and Praise)."; Noor, "Muqaddimah of Ibn Al-salah.".} \textit{Takhrij} kemudian ditransformasikan menjadi ekstraksi fitur komputasional plus pemetaan graf semantik; sedangkan \textit{jarh-ta'dil}, rintisan 'Ali b. al-Madini yang diwarisi Ibn al-Salah--Mamluk, disahkan lewat kajian jejaring atas \textit{Sahih Bukhari} berpendekatan jaringan.\footnote{Wasman, "A Critical Approach to Prophetic Traditions."; A. B. Kamran et al., "SemanticHadith: An ontology-driven knowledge graph for the hadith corpus," \\textit{Journal of Web Semantics}, vol. 78, hlm. 100797, 2023; Su, "'ali B. Al-madini (161-234/778-849)."; T. Alam and J. Schneider, "Social network analysis of hadith narrators from Sahih Bukhari," \\textit{Proceedings of 2020 7th IEEE International Conference on Behavioural and Social Computing (BESC)}, hlm. 1--6, 2020.} Arsip \textit{maudu'at} rintisan Ibn al-Jauzi hingga al-Albani dialihfungsikan menjadi materi latih algoritme deteksi; telaah Ibn Kasir atas \textit{Sahihayn} menunjukkan tiada wilayah luput dari penilaian.\footnote{Noor, "Muqaddimah of Ibn Al-salah."; Calgan, "Analysis of Ibn Kathir's Content.".} Separuh akhir rumusan ketiga pun menyimpulkan: kekebalan epistemologis hadis disangga oleh kesediaan diuji ulang lewat mekanisme kelembagaan yang ajek antar-generasi.

Kelima, makalah ini merangkai telaah historis-filologis klasik dengan pendekatan komputasional-digital --- meramu benang periodisasi, transmisi, dan pemalsuan secara bersintesis.\footnote{M. A. Abdulrahman, "The future of hadith studies in the digital age: Opportunities and challenges," \\textit{Journal of Ecohumanism}, vol. 3, no. 8, 2024; K. H. Bin Jamil, "Between traditionalising and futuring: Applying the broader maqasid paradigm to hadith studies," \\textit{Asia Pacific Journal of Educators and Education}, vol. 39, no. 2, hlm. 183--196, 2024.} Peranti \textit{pretrained}, \textit{AraBERT}, dan graf-\textit{transformer} diposisikan meneruskan \textit{tathabbut}, \textit{rihlah}, \textit{jarh-ta'dil} --- bukan menggantikan ulama --- sejalan dengan kerangka \textit{maqasid} sebagaimana dirumuskan.\footnote{Bin Jamil, "Between Traditionalising and Futuring."; Abdulrahman, "The Future of Hadith Studies in the Digital Age.".} Urgensinya berpangkal pada konten viral yang hampa validitas dan \textit{ChatGPT} sebagai rujukan belajar, sehingga mendesak pengalihdigitalan \textit{takhrij} serta literasi hadis yang kritis-klasik.\footnote{Z. Abdullah et al., "Viralitas dan validitas: Digitalisasi takhrij hadis di era media sosial," \\textit{Dirayah: Jurnal Ilmu Hadis}, vol. 6, no. 2, hlm. 381--390, 2026; H. Syukur et al., "Living hadith in the age of artificial intelligence: A phenomenological study of the use of ChatGPT as a source for studying hadith among college students," \\textit{Asian Journal of Islamic Studies and Da'wah}, vol. 4, no. 3, hlm. 189--204, 2026.} Karenanya ilmu hadis tampil sebagai \textit{living hadith} yang lentur --- membentang dari \textit{sahifah} klasik ke pangkalan data masa kini.\footnote{Abdulrahman, "The Future of Hadith Studies in the Digital Age."; A. M. B. Matin Bin Salman, "Reconstructing hadith discourse in the digital age: From text to discourse," \\textit{Journal of Ecohumanism}, vol. 4, no. 1, hlm. 1--11, 2025.}

\phantomsection
\subsection*{B. Saran}
\addcontentsline{toc}{subsection}{B. Saran}

\begin{enumerate}[leftmargin=1.27cm,itemsep=2pt,parsep=2pt]
\item Perbandingan lintas mazhab-antargenerasi soal tafsir larangan penulisan layak diprioritaskan, sebab rekonsiliasi kronologis disandarkan\footnote{Lutfianto, "History of Hadith Writing."; Watukila, "Historiografi Hadis Di Masa Mutaqaddimin.".} pada kepustakaan Sunni, 'Ajjaj al-Khatib, dan telaah Azami. Kepustakaan Syiah yang diperhadapkan dengan riset orientalis akan menguji daya laku umum kesepakatan \textit{nasakh} pelarangan lewat pembolehan penulisan menyeluruh.

\item Tolok ukur Arab-Indonesia untuk pemeriksaan digital memadukan petanda \textit{sanad} (kesinambungan dan sentralitas jaringan) dengan petanda \textit{matan} (gaya, sentimen, koherensi) dalam penilaian terbuka hasil kesepakatan.\footnote{E. Mohamed and R. Sarwar, "Linguistic features evaluation for hadith authenticity through automatic machine learning," \\textit{Digital Scholarship in the Humanities}, vol. 37, no. 3, hlm. 830--843, 2022; J. Alghamdi et al., "Pretrained models against traditional machine learning for detecting fake hadith," \\textit{Electronics}, vol. 14, no. 17, hlm. 3484, 2025.} Berikutnya, korpus \textit{maudu'at} beranotasi \textit{'illah} dialihbentuk secara digital sebagai bahan latih penimbang setara model \textit{pretrained} dan \textit{AraBERT} secara berdampingan.\footnote{M. Shaaban et al., "Prophetic authorial style modeling for detecting fabricated hadiths using AraBERT," \\textit{Journal of Computing & Biomedical Informatics}, vol. 10, no. 2, 2026; T. Tarmom et al., "Deep learning vs compression-based vs traditional machine learning classifiers to detect hadith authenticity," \\textit{Communications in Computer and Information Science}, hlm. 206--222, 2022.}

\item Literasi keagamaan kaum muda seyogianya membekalkan \textit{takhrij} digital secara bertahap --- dari penelusuran daring sampai pengecekan \textit{sanad}-\textit{matan} --- untuk menanggapi \textit{ChatGPT} sebagai bahan belajar dan viralnya hadis tanpa keabsahan.\footnote{Syukur, "Living Hadith in the Age of Artificial Intelligence."; Abdullah, "Viralitas Dan Validitas.".} Riset aksi (\textit{action research}) literasi hadis digital yang telah dirintis selayaknya diperbanyak di jenjang menengah dan dakwah dalam jaringan.\footnote{T. Supriyadi et al., "Action research in hadith literacy: A reflection of hadith learning in the digital age," \\textit{International Journal of Learning, Teaching and Educational Research}, vol. 19, no. 5, hlm. 99--124, 2020; Abdulrahman, "The Future of Hadith Studies in the Digital Age.".}

\item Kritik \textit{sanad}-\textit{matan} dalam kurikulum keislaman --- dari pesantren hingga kampus --- seyogianya mengangkat \textit{jarh-ta'dil}, \textit{takhrij}, dan telaah kontekstual menjadi kecakapan yang melampaui hafalan semata.\footnote{Muhamed Ali, "Al-jarh Wa Al-ta'dil (criticism and Praise)."; Wasman, "A Critical Approach to Prophetic Traditions.".} Kebangkitan Mamluk lewat mediasi \textit{Muqaddimah} Ibn al-Salah menunjukkan bahwa pembakuan kurikulum selalu lahir untuk menjawab kegentingan; teladan ini pantas ditempuh ulang saat ini.\footnote{Noor, "Muqaddimah of Ibn Al-salah."; Bin Jamil, "Between Traditionalising and Futuring.".}

\item Kajian prosopografis jejaring periwayat Hijaz, Irak, Syam, dan Mesir mengawinkan telaah jejaring dengan graf pengetahuan berbasis \textit{transformer} untuk membidik anomali transmisi berskala luas.\footnote{Alam, "Social Network Analysis of Hadith."; M. A. Mosa, "Synergizing structure and semantics: A knowledge graph-transformer framework for narrator disambiguation in hadith networks," \\textit{Digital Scholarship in the Humanities}, 2025.} Lewat pendekatan itu, ketaksaan identitas narator dapat disempurnakan, dan dugaan kehati-hatian berlebih pada wilayah tertentu dapat diuji memakai data terbuka-tereplikasi menurut telaah.\footnote{Kamran, "Semantichadith."; Su, "'ali B. Al-madini (161-234/778-849).".}

\end{enumerate}
